\documentclass[rnd]{mas_proposal}
\usepackage[utf8]{inputenc}

\usepackage[margin=0.7in]{geometry}
\usepackage{tikz}
\usepackage{float}
\usepackage{amsmath}
\usepackage{amsfonts}
\usepackage{amssymb}
\usepackage{graphicx}
\usepackage[numbers]{natbib}
\usepackage{tcolorbox}
\usepackage{listings}
\usepackage{xcolor}

% Code style
\lstset{
basicstyle=\ttfamily\footnotesize,
backgroundcolor=\color{gray!10},
keywordstyle=\color{blue},
commentstyle=\color{green!50!black},
stringstyle=\color{orange},
breaklines=true,
frame=single,
showstringspaces=false
}

\title{Progress Report:Rich Descriptive Models for Reusable ROS Software Components}
\author{Riddhesh Pradeep More}
\supervisors{Prof. Sebastian Houben \ Dr. Bj"orn Kahl}
\date{October 2025}

\begin{document}

\maketitle
\pagestyle{plain}

% ---------------- EXECUTIVE SUMMARY ----------------
\section{Executive Summary}
This report details the progress of the R\&D project, titled "Rich Descriptive Models for Reusable ROS Software Components." The project has successfully pivoted from a purely Model-Driven Engineering (MDE) approach to a more pragmatic and highly functional hybrid system. This change was driven by a need to create a solution that is not only technically rigorous but also intuitive and scalable for the end user.

The core of this lies in the integration of semantic web technologies with modern AI, specifically \textbf{vector embeddings}. This has enabled in delivering a functional, multi-layered prototype well ahead of its initial schedule. The current system provides a semantic search and discovery interface for ros components, demonstrating a significant advancement over the original proposal's scope and a substantial contribution to the IAM4RAIL project.

% ---------------- ORIGINAL IDEA (NEW CHAPTER 2) ----------------
\section{Original Idea: Addressing the Semantic Gap}
The genesis of this R&D project was a critical and persistent bottleneck in modern robotics development, often referred to as the "semantic gap" in the Robot Operating System (ROS) ecosystem \citep{quigley2009ros}. While ROS has become a foundational framework, providing a modular architecture and a comprehensive ecosystem of packages, it fundamentally lacks a standardized, machine-interpretable semantic layer. This means that while the framework defines standards for communication and node interaction, it does not have a formalized way to describe the function, compatibility, or constraints of its software and hardware components.

As a result, developers are forced to rely heavily on ad hoc, human-readable documentation and keyword searches, which are frequently incomplete or outdated. This reliance on manual exploration significantly impedes component discovery and reuse, thereby hindering the scalability of robotic solutions across different domains. The project was initially conceived to directly address this fundamental deficiency by providing a structured, verifiable approach to component description.

\subsection{Originally Envisioned Approach}
The project was originally conceived to follow a classic Model-Driven Engineering (MDE) methodology to create formal, machine-interpretable models for ROS components \citep{hammoudeh2019bootstrapping, zander2016model, brugali2015model}. The plan involved several distinct phases.

First, a \textbf{System Analysis} phase was planned to examine component configurations on existing outdoor robotic platforms used for inspection and maintenance. This work would be done in collaboration with robotics teams to identify and classify functional groupings such as SLAM, path planning, and actuation.

Based on this analysis, the next step was \textbf{Ontology and Taxonomy Design}. This involved designing a formal hierarchy of component classes (e.g., \texttt{LocalizationNode}, \texttt{SensorDriver}) and defining their functional and non-functional properties—such as message types and hardware compatibility—using a standard format like OWL or RDF.

Finally, the plan was to implement these models by building a \textbf{Searchable Component Database}. This involved creating a semantic repository using an RDF store like Blazegraph and developing a user interface with a framework like NiceGUI for querying and visualizing the component data. The overarching goal was to produce a structured, searchable knowledge base that would enhance component discovery and integration efficiency.

% ---------------- IMPLEMENTATION (NEW CHAPTER 3) ----------------
\section{Implementation and Progress of the Hybrid System}
To address the practical challenges of the MDE approach, particularly its complexity and rigidity for developers, the project pivoted towards a more flexible hybrid architecture. This new direction combines the strengths of formal semantic models with the intuitive power of modern AI-driven search. The goal was to create a system that allows developers to find components using natural language queries, bridging the gap between user intent and structured component data.

\subsection{System Architecture}
The current architecture integrates structured knowledge with semantic meaning (see Fig.~\ref{fig:architecture}).

\begin{figure}[H]
\centering
% The TikZ code has been replaced with this command to include the PNG.
% Make sure your image file is named "system_architecture.png" and is
% in the same folder as your.tex file.
\includegraphics[width=\textwidth]{images/System_Arch.png}
\caption{High-level system architecture}
\label{fig:architecture}
\end{figure}

\subsection{Pivot from SPARQL to Solr}
Originally, the system relied on SPARQL queries against an RDF store (Blazegraph).
\textbf{Challenges:}
\begin{itemize}
\item Developers found SPARQL unintuitive and keyword brittle.
\item Limited support for semantic synonyms or intent-based queries.
\item Performance issues on larger graphs.
\end{itemize}

\textbf{Change:} Migrated to \textbf{Apache Solr 9.x} with hybrid lexical + vector search.
\textbf{Why:} Solr natively supports dense vector fields and hybrid scoring, enabling natural language queries while retaining ontology precision.

\subsection{Integration of Vector Embeddings}
\begin{itemize}
\item Implemented \textbf{all-MiniLM-L6-v2} \citep{wang2020minilm} (384D) embeddings for ROS package descriptions and topics.
\item User queries are vectorized for k-NN semantic similarity.
\item Embeddings bridge the gap between developer phrasing and component metadata, leveraging models like Sentence-BERT \citep{reimers2019sentencebert}.
\end{itemize}

% ---------------- MATHEMATICAL FORMULATION ----------------
\section{Mathematical Formulation of Semantic Similarity}

The core innovation of this hybrid system lies in its mathematical approach to semantic similarity calculation, which bridges the gap between human-interpretable natural language queries and machine-processable vector representations. This section presents the formal mathematical foundation underlying the semantic search capabilities.

\subsection{Vector Embedding Generation}

The system employs Sentence-BERT \citep{reimers2019sentencebert}, specifically the \texttt{all-MiniLM-L6-v2} model, to transform textual component descriptions into dense vector representations. For a given component description text $T$, the embedding process follows:

\begin{align}
\mathbf{E} &= \text{BERT}(\text{tokenize}(T)) \in \mathbb{R}^{L \times 384} \label{eq:bert_encoding}\\
\mathbf{v}_{\text{raw}} &= \frac{1}{L} \sum_{i=1}^{L} \mathbf{E}_i \label{eq:mean_pooling}\\
\mathbf{v} &= \frac{\mathbf{v}_{\text{raw}}}{\|\mathbf{v}_{\text{raw}}\|_2} \label{eq:l2_normalization}
\end{align}

where $L$ represents the sequence length, $\mathbf{E}_i$ denotes the $i$-th token embedding, and $\mathbf{v} \in \mathbb{R}^{384}$ is the final L2-normalized component vector with $\|\mathbf{v}\|_2 = 1$.

\subsection{Component Text Synthesis}

Before embedding generation, component metadata is synthesized into a comprehensive textual representation. For a component $c$ with attributes $\{name, type, description, package, topics\}$, the synthesis function is defined as:

\begin{equation}
T(c) = \text{``Component name: ''} \oplus c.name \oplus \text{`` Type: ''} \oplus c.type \oplus \cdots
\end{equation}

where $\oplus$ denotes string concatenation, ensuring rich semantic context for embedding generation.

\subsection{Cosine Similarity Calculation}

The fundamental similarity measure between components is cosine similarity. Given two normalized vectors $\mathbf{v}_1, \mathbf{v}_2 \in \mathbb{R}^{384}$ with $\|\mathbf{v}_1\|_2 = \|\mathbf{v}_2\|_2 = 1$, the cosine similarity simplifies to:

\begin{equation}
\text{sim}(\mathbf{v}_1, \mathbf{v}_2) = \cos(\theta) = \mathbf{v}_1 \cdot \mathbf{v}_2 = \sum_{i=1}^{384} v_{1i} \cdot v_{2i} \label{eq:cosine_similarity}
\end{equation}

where $\theta$ is the angle between vectors, and the result $\text{sim}(\mathbf{v}_1, \mathbf{v}_2) \in [-1, 1]$ with $1$ indicating identical semantic meaning and $-1$ indicating opposite meanings.

\subsection{k-Nearest Neighbors Search}

For a query vector $\mathbf{q}$ and a database of $n$ component vectors $\{\mathbf{c}_1, \mathbf{c}_2, \ldots, \mathbf{c}_n\}$, the k-NN search algorithm computes:

\begin{align}
S &= \{(i, \text{sim}(\mathbf{q}, \mathbf{c}_i)) : i = 1, 2, \ldots, n\} \label{eq:similarity_set}\\
S_{\text{sorted}} &= \text{sort}(S, \text{key}=\text{similarity}, \text{desc}=\text{True}) \label{eq:sorting}\\
\text{Result} &= S_{\text{sorted}}[1:k] \label{eq:top_k}
\end{align}

The computational complexity is $O(n \cdot d)$ where $d = 384$ is the vector dimension, making it linear in the number of components.

\subsection{Hybrid Search Scoring}

The system combines traditional text-based search scores with vector similarity through a weighted linear combination. Let $s_{\text{text}} \in [0, 1]$ be the normalized BM25 relevance score from Apache Solr, and $s_{\text{vector}} \in [-1, 1]$ be the cosine similarity score. The hybrid score is computed as:

\begin{align}
s_{\text{vector,norm}} &= \frac{s_{\text{vector}} + 1}{2} \in [0, 1] \label{eq:vector_normalization}\\
s_{\text{hybrid}} &= \alpha \cdot s_{\text{vector,norm}} + (1 - \alpha) \cdot s_{\text{text}} \label{eq:hybrid_score}
\end{align}

where $\alpha \in [0, 1]$ is the semantic weight parameter (default $\alpha = 0.7$), emphasizing vector similarity while preserving text relevance signals.

\subsection{Distance Metrics and Alternatives}

While cosine similarity is the primary metric, the framework supports alternative distance measures:

\begin{align}
\text{Euclidean Distance:} \quad d_E(\mathbf{v}_1, \mathbf{v}_2) &= \|\mathbf{v}_1 - \mathbf{v}_2\|_2 \label{eq:euclidean}\\
\text{Manhattan Distance:} \quad d_M(\mathbf{v}_1, \mathbf{v}_2) &= \sum_{i=1}^{384} |v_{1i} - v_{2i}| \label{eq:manhattan}\\
\text{Angular Distance:} \quad d_A(\mathbf{v}_1, \mathbf{v}_2) &= \frac{\arccos(\mathbf{v}_1 \cdot \mathbf{v}_2)}{\pi} \label{eq:angular}
\end{align}

Each metric can be converted to a similarity score using $\text{sim} = \frac{1}{1 + d}$ for distance-based measures.

\subsection{Mathematical Properties and Guarantees}

The vector space exhibits several important mathematical properties:

\begin{itemize}
\item \textbf{Metric Properties}: Cosine similarity satisfies symmetry ($\text{sim}(\mathbf{u}, \mathbf{v}) = \text{sim}(\mathbf{v}, \mathbf{u})$) and bounded range ($[-1, 1]$).
\item \textbf{Dimensionality}: Fixed 384-dimensional space ensures consistent comparison across all components.
\item \textbf{Normalization}: L2 normalization ensures similarity is invariant to vector magnitude, focusing purely on directional similarity.
\item \textbf{Semantic Preservation}: The Sentence-BERT model preserves semantic relationships in the vector space, enabling meaningful similarity calculations.
\end{itemize}

This mathematical foundation provides a rigorous basis for semantic component discovery, enabling queries that go beyond keyword matching to capture semantic intent and functional similarity.

% ---------------- KNOWLEDGE REPRESENTATION ----------------
\section{Knowledge Representation: The Ontology and Triple Model}
The approach of using a formal model to describe ROS packages has strong scientific precedent. Research in this area aims to solve the exact "semantic gap" this project targets, for example, by implementing knowledge bases for reusable ROS components. The key takeaway from prior research is the direct validation of using a knowledge graph—built with an ontology and populated with structured data—as a practical and effective method for enabling semantic queries for component discovery. This validates our choice of technologies and methodology.

\subsection{Building the Foundation: The Ontology Model}
Before facts about ROS components can be stored, a structural blueprint is required. This is the role of the \textbf{ontology}, a formal specification of concepts and their relationships \citep{gruber1993translation}. Defined in a Turtle file (\texttt{ros_package_ontology.ttl}), it acts as the schema for our knowledge graph. The ontology is built using standard W3C recommendations: RDF Schema (RDFS) for class hierarchies and Web Ontology Language (OWL) for defining complex properties.

\begin{tcolorbox}[colback=blue!5!white,colframe=blue!75!black,title=Key Ontology Constructs]
\begin{enumerate}
\item \textbf{Defining Classes (The "Nouns"):} A hierarchy of component types was created. This allows the system to infer, for example, that a \texttt{LocalizationPackage} is a specific type of \texttt{AlgorithmPackage}.
\begin{lstlisting}[language=bash]

A LocalizationPackage is a type of AlgorithmPackage

ros:LocalizationPackage rdfs:subClassOf ros:AlgorithmPackage.
\end{lstlisting}

\item \textbf{Defining Properties (The "Verbs"):} The relationships components can have are defined, such as dependencies, hardware requirements, and compatibility.

\begin{lstlisting}[language=bash]

A property to link a package to required hardware

ros:requiresHardware a owl:ObjectProperty.
\end{lstlisting}
\end{enumerate}
\end{tcolorbox}
This ontology provides the essential structure for our knowledge, ensuring data is consistent and queryable.

\subsection{Populating the Knowledge Graph with Triples}
With the ontology as our blueprint, we populate it with facts about real ROS packages. This process turns unstructured documentation into structured, machine-readable knowledge based on the Resource Description Framework (RDF) model. Each piece of information is stored as a \textbf{Subject-Predicate-Object triple}, which forms the atomic unit of the knowledge graph.

\begin{tcolorbox}
The human-readable fact, "The AMCL package is a localization algorithm that requires a laser scanner," is translated into the following triples:
\begin{lstlisting}[language=bash]
@prefix ros: http://example.org/ros/.
@prefix hw: http://example.org/hardware/.
@prefix rdf: http://www.w3.org/1999/02/22-rdf-syntax-ns#.
@prefix rdfs: http://www.w3.org/2000/01/rdf-schema#.

Subject Predicate Object

#----------------   -------------------   ------------------------
ros:amcl            rdf:type               ros:LocalizationPackage ;
rdfs:label             "AMCL" ;
ros:requiresHardware   hw:LaserScanner ;
ros:compatibleWith     ros:move_base.
\end{lstlisting}
By repeating this process, we build a complex graph, of interconnected knowledge that the system can navigate to answer complex queries. This structured data is then indexed in Apache Solr to become a searchable knowledge base.
\end{tcolorbox}

\subsection{Knowledge Base Improvements}
The RDF knowledge base was expanded with well-known ROS packages for mobile robots and their underlying algorithms:
\begin{itemize}
\item Localization: \texttt{AMCL}, \texttt{robot_localization}, \texttt{ORB-SLAM2}.
\item Navigation: \texttt{Navigation2}, \texttt{TEB Planner}.
\item Perception: \texttt{darknet_ros} (based on YOLO), \texttt{apriltag_ros}.
\end{itemize}
This structured data forms a knowledge graph, whose schema and population are detailed below.

% ---------------- RESULTS ----------------
\section{Results and Evaluation}
\begin{itemize}
\item Hybrid search pipeline operational: user queries can be matched via both BM25 text similarity and embedding similarity.
\item Prototype supports ∼20 ROS components with structured metadata.
\item Semantic queries such as \textit{``Find components that localize a mobile robot''} successfully retrieve \texttt{AMCL} and \texttt{robot_localization}.
\item Limitation: The models are manually coded.
\end{itemize}

\subsection{Automated Knowledge Graph Construction: TO DO}
To overcome manual TTL authoring, a prototype \textbf{LangChain agent pipeline} can be introduced. This aligns with recent work in using LLMs for automated knowledge graph construction \citep{geng2024survey}.
\begin{enumerate}
\item Extract metadata from \texttt{package.xml} and \texttt{README.md}.
\item Generate Subject-Predicate-Object triples guided by ontology schema.
\item Validate triples via \texttt{rdflib} and index in Solr.
\end{enumerate}

This agent-based approach allows scalable ingestion of ROS repositories.

% ---------------- FUTURE WORK ----------------
\section{Future Work}
\begin{itemize}
\item Automate ingestion of thousands of ROS packages via LLM agent.
\item Add dependency visualization (topics, messages, versions).
\item Conduct developer user studies to validate semantic relevance.
\end{itemize}

% ---------------- CONCLUSION ----------------
\section{Conclusion}
The project has successfully transitioned from a brittle, formal-only approach to a scalable hybrid semantic search architecture. By combining a structured knowledge base (RDF), vector embeddings, and a powerful search backend (Solr), the system demonstrates a practical and effective method for allowing robotics developers to query for components based on meaning and intent, rather than just keywords.

This progress establishes a strong foundation for future scaling and automated knowledge ingestion, marking a significant step forward in addressing the semantic gap in modern robotics ecosystems.

\begin{thebibliography}{99}

\bibitem{brugali2015model}
Brugali, D. (2015). Model-driven software engineering in robotics: Models are designed to use the relevant things, thereby reducing the complexity and cost in the field of robotics. \textit{IEEE Robotics Automation Magazine}, 22(3), 155–166.

\bibitem{bruyninckx2013brics}
Bruyninckx, H., Klotzbücher, M., Hochgeschwender, N., Kraetzschmar, G., Gherardi, L., & Brugali, D. (2013). The brics component model: a model-based development paradigm for complex robotics software systems. In \textit{Proceedings of the 28th Annual ACM Symposium on Applied Computing (SAC '13)} (pp. 1758–1764). ACM.

\bibitem[Ciccozzi et al.(2022)]{ciccozzi2022model}
Ciccozzi, F., et al. (2022). Model-driven engineering for mobile robotic systems: A systematic mapping study. \textit{Software and Systems Modeling}, 21, 19–49.

\bibitem[de Lara et al.(2022)]{kolovos2022lowcode}
de Lara, J., et al. (2022). Low-code development and model-driven engineering: Two sides of the same coin? \textit{Software and Systems Modeling}, 21, 437–446.

\bibitem{durschmid2024rosinfer}
Dürschmid, T., Timperley, C. S., Garlan, D., & Le Goues, C. (2024). Rosinfer: Statically inferring behavioral component models for ros-based robotics systems. In \textit{Proceedings of the IEEE/ACM 46th International Conference on Software Engineering (ICSE '24)}. ACM.

\bibitem[Geng et al.(2024)]{geng2024survey}
Geng, Y., Li, Y., Zhang, W., Wang, P., Liu, Y., & Mei, J. (2024). A Survey on Large Language Models for Aided Knowledge Graph Construction. \textit{arXiv preprint arXiv:2402.08399}.

\bibitem[Gruber(1993)]{gruber1993translation}
Gruber, T. R. (1993). A translation approach to portable ontology specifications. \textit{Knowledge acquisition}, 5(2), 199-220.

\bibitem[Hammoudeh Garcia et al.(2019)]{hammoudeh2019bootstrapping}
Hammoudeh Garcia, N., Lüdtke, M., Kortik, S., Kahl, B., & Bordignon, M. (2019). Bootstrapping mde development from ros manual code-part 1: Metamodeling. In \textit{2019 IEEE International Conference on Robotic Computing (IRC)}.

\bibitem[Hammoudeh Garcı́a et al.(2021)]{hammoudeh2021bootstrapping}
Hammoudeh Garcı́a, N., et al. (2021). Bootstrapping mde development from ros manual code: Part 2—model generation and leveraging models at runtime. \textit{Software and Systems Modeling}, 20, 2047–2070.

\bibitem{iam4rail2025}
IAM4RAIL Project. (2025). Iam4rail: European railway research project. Retrieved from https://projects.rail-research.europa.eu/eurail-fp3/

\bibitem[Neto et al.(2020)]{neto2020applying}
Neto, T., Arrais, R., Sousa, A., Santos, A., & Veiga, G. (2020). Applying Software Static Analysis to ROS: The Case Study of the FASTEN European Project. In \textit{Robot 2019: Fourth Iberian Robotics Conference} (pp. 632–644). Springer.

\bibitem[Quigley et al.(2009)]{quigley2009ros}
Quigley, M., Conley, K., Gerkey, B., Faust, J., Foote, T., Leibs, J.,... & Ng, A. Y. (2009). ROS: an open-source Robot Operating System. In \textit{ICRA workshop on open source software} (Vol. 3, No. 3.2).

\bibitem{reimers2019sentencebert}
Reimers, N., & Gurevych, I. (2019). Sentence-BERT: Sentence Embeddings using Siamese BERT-Networks. In \textit{Proceedings of the 2019 Conference on Empirical Methods in Natural Language Processing}.

\bibitem{scott2025ros}
Scott, K., & Foote, T. (2025). 2024 ros metrics report. Open Robotics Foundation. Retrieved from https://www.ros.org/news/2025/01/2024-ros-metrics-report.html

\bibitem{trezzy2021applying}
Trezzy, M., Ober, I., Ober, I., & Oliveira, R. (2021). Applying mde to ros systems: A comparative analysis. \textit{Scientific Annals of Computer Science}, 31, 111–144.

\bibitem{wang2020minilm}
Wang, W., Wei, F., Dong, L., Bao, H., Yang, N., & Zhou, M. (2020). MiniLM: Deep Self-Attention Distillation for Task-Agnostic Compression of Pre-Trained Transformers. \textit{arXiv preprint arXiv:2002.10957}.

\bibitem[Zander et al.(2016)]{zander2016model}
Zander, S., Heppner, G., Neugschwandtner, G., Awad, R., Essinger, M., & Ahmed, N. (2016). A model-driven engineering approach for ros using ontological semantics. \textit{arXiv preprint arXiv:1601.03998}.

\end{thebibliography}

\end{document}